\documentclass[10pt,a4paper]{amsart}
\usepackage[margin=19mm]{geometry}
\usepackage{amsmath,amssymb,mathtools}
\usepackage[scr=rsfs,cal=euler]{mathalfa}
\usepackage{xfrac}
\usepackage[unicode,hidelinks,bookmarks=false]{hyperref}
\newcommand{\ie}{i.e.\ }
\newcommand{\cf}{cf.\ }
\newcommand{\resp}{respectively\ }
\newcommand{\Subsection}{\subsection}
\providecommand{\nobreakdash}{\nobreak\hbox{-}}
\setlength{\emergencystretch}{2em}
\multlinegap0pt
\usepackage{amssymb}
\newtheorem{proposition}{Proposition}
\newtheorem{theorem}{Theorem}
\newcommand{\Qdot}{\dot{Q\omega}}
\newcommand{\R}{\mathbb{R}}
\newcommand{\barE}{\bar{E}}
\newcommand{\barp}{\bar{p}}
\newcommand{\barq}{\bar{q}}
\newcommand{\barrho}{\bar{\rho}}
\newcommand{\bartau}{\bar{\tau}}
\newcommand{\sgs}{\mathrm{sgs}}
\newcommand{\tilu}{\widetilde{u}}
\title{\bfseries Diffeomorphism Invariance of the Favre-Filtered Compressible Navier--Stokes System under Spacetime Dilation}
\author{Yuanya Li}
\date{Source: arXiv:2609.11566v1 (CC BY 4.0)}
\begin{document}
\maketitle

\noindent\emph{Source and licence.} \bfseries Diffeomorphism Invariance of the Favre-Filtered Compressible Navier--Stokes System under Spacetime Dilation. Version arXiv:2609.11566v1 (CC BY 4.0), distributed by arXiv under the Creative Commons Attribution 4.0 International licence, http://creativecommons.org/licenses/by/4.0/, as recorded in the arXiv licence metadata for this exact version and shown on the abstract page https://arxiv.org/abs/2609.11566v1. Full licence notice: \texttt{licenses/arxiv-2609.11566-CC-BY-4.0.txt}.\par\smallskip

\noindent\textit{Editorial scope.} Counted unit: the article's theorem thm:equivariance (source0.tex lines 741--754). The article's complete original proof of it is delivered with it (source0.tex lines 756--803, one \texttt{proof} environment of 48 lines). No mathematics is written, repaired or re-derived: the statement and the proof are the article's own lines, in the article's own notation, and no retained line is substituted or re-flowed.  Retained uncounted as the source-local context the proof needs: lines 534--540; lines 589--595; lines 656--689.  Nothing was omitted from any retained span, so the package carries no omission record.  No retained line contains a citation, so the wrapper carries no bibliography and no bibliography entry.  No retained line contains a figure environment, an image-inclusion command or drawing code, so no image is delivered and the package declares no asset dependency.  Editorial formatting only, in addition: an AMS \texttt{amsart} preamble that restates the article's own declarations and macros used by the retained lines, namely the environments \texttt{proposition}, \texttt{theorem} on separate counters, together with the standard packages those lines need.  The retained environments print in this document as Proposition 1, Theorem 1 in order of appearance, whereas the article prints the same items with the numbering of its own full document.  The complete native source of the article as distributed in the arXiv e-print for this exact version is bundled unchanged as \texttt{sources/210033/source0.tex}.
\begin{proposition}
\label{prop:continuity}
\[
R_0[\Phi_L(u)](t,x)=L\,R_0[u](Lt,Lx).
\]
Consequently, $R_0[\Phi_L(u)]=0$ if and only if $R_0[u]=0$.
\end{proposition}

\begin{proposition}
\label{prop:momentum}
\[
R_i[\Phi_L(u)](t,x)=L\,R_i[u](Lt,Lx).
\]
Hence $R_i[\Phi_L(u)]=0$ iff $R_i[u]=0$.
\end{proposition}

\begin{proposition}
\label{prop:energy}
\[
R_2[\Phi_L(u)](t,x)
=L\Bigl[\partial_T(\barrho\barE)
+\partial_{X_i}\bigl((\barrho\barE+\barp)\tilu_i+\barq_i
-\tilu_j\bartau_{ij}\bigr)
+\partial_{X_i}H^{\sgs}_i
+\partial_{X_i}\sigma^{\sgs}_{ij}\Bigr](Lt,Lx)
+\Qdot(Lt,Lx).
\]
Equivalently, if we define the \emph{scaled energy residual}
\[
\widehat R_2[u]:=\frac{1}{L}R_2[u]\quad\text{in dilated coordinates},
\]
then in $(T,X)$ coordinates the equation $R_2[u]=0$ becomes
\[
\frac{\partial(\barrho\barE)}{\partial T}
+\frac{\partial}{\partial X_i}\bigl[(\barrho\barE+\barp)\tilu_i
+\barq_i-\tilu_j\bartau_{ij}\bigr]
+\frac{\partial H^{\sgs}_i}{\partial X_i}
+\frac{\partial\sigma^{\sgs}_{ij}}{\partial X_i}
=-\Qdot,
\]
which is exactly
\[
\frac{\partial(\barrho\barE)}{\partial(Lt)}
+\frac{\partial}{\partial(Lx_i)}\bigl[(\barrho\barE+\barp)\tilu_i
+\barq_i-\tilu_j\bartau_{ij}\bigr]
+\frac{\partial H^{\sgs}_i}{\partial(Lx_i)}
+\frac{\partial\sigma^{\sgs}_{ij}}{\partial(Lx_i)}
=-\Qdot.
\]
\end{proposition}

\begin{theorem}[Equivariance of the residual map]
\label{thm:equivariance}
Let $\mathcal P=(R_0,R_1,R_2,R_3,R_4):\mathcal C\to C^\infty(M,\R^5)$ be the
full residual map (one continuity + three momentum components + one energy).
Then for every $L>0$ and $u\in\mathcal C$,
\[
\mathcal P(\Phi_L(u))=L\,\phi_L^*(\mathcal P(u))
\]
componentwise, where $\phi_L^*$ acts on the $\R^5$-valued function
$\mathcal P(u)$ by $(\phi_L^*F)(t,x)=F(Lt,Lx)$.  In particular,
\[
\mathcal P(\Phi_L(u))=0\quad\Longleftrightarrow\quad\mathcal P(u)=0.
\]
\end{theorem}

\begin{proof}
The continuity and momentum components were shown in
Propositions \ref{prop:continuity} and \ref{prop:momentum} to satisfy
$R_\alpha[\Phi_L(u)](t,x)=L\,R_\alpha[u](Lt,Lx)$ for
$\alpha=0,1,2,3$.  For the energy component, \ref{prop:energy} shows that in
the passive coordinate representation the equation is identical in form; in
the active representation,
\[
R_4[\Phi_L(u)](t,x)
=L\bigl[\text{derivative terms of }R_4[u]\bigr](Lt,Lx)+\Qdot(Lt,Lx).
\]
Now observe that the original energy equation is $R_4[u]=0$, i.e.
\[
\bigl[\text{derivative terms}\bigr](t,x)+\Qdot(t,x)=0.
\]
Evaluating at $(Lt,Lx)$:
\[
\bigl[\text{derivative terms}\bigr](Lt,Lx)+\Qdot(Lt,Lx)=0.
\]
Multiplying by $L$:
\[
L\bigl[\text{derivative terms}\bigr](Lt,Lx)+L\,\Qdot(Lt,Lx)=0.
\]
This differs from $R_4[\Phi_L(u)]$ by the source term: $R_4[\Phi_L(u)]$ has
$\Qdot(Lt,Lx)$ rather than $L\,\Qdot(Lt,Lx)$.  Therefore the strict
equivariance $R_4[\Phi_L(u)]=L\,\phi_L^*R_4[u]$ holds if and only if we
regard the source as part of the field tuple and the equation as
\emph{inhomogeneous}.  However, the \emph{solution-set preservation} still
holds:

($\Rightarrow$) If $\mathcal P(\Phi_L(u))=0$, then in particular
$R_4[\Phi_L(u)](t,x)=0$ for all $(t,x)$:
\[
L\,D[u](Lt,Lx)+\Qdot(Lt,Lx)=0,
\]
where $D[u]$ denotes the derivative part.  Since $\phi_L$ is surjective, as
$(t,x)$ ranges over $M$, $(Lt,Lx)$ ranges over $M$, so
\[
L\,D[u](p)+\Qdot(p)=0\quad\forall p\in M.
\]
This is \emph{not} the original equation $D[u]+\Qdot=0$ unless $L=1$.

We must therefore refine the statement.  The correct transformation of an
inhomogeneous PDE under a diffeomorphism requires the source to transform as
a density of appropriate weight, or equivalently we must include the source
in the field tuple and transform it as part of the section.  We now do this
properly.
\end{proof}

\end{document}
